% v6: supervisor-review draft after annotated v5 feedback
\documentclass[conference,letterpaper]{IEEEtran}
\IEEEoverridecommandlockouts
% The preceding line is only needed to identify funding in the first footnote. If that is unneeded, please comment it out.
\usepackage{cite}
\usepackage{amsmath,amssymb,amsfonts}
\usepackage{algorithm}
\usepackage{algpseudocode}
\usepackage{graphicx}
\usepackage{subcaption}
\usepackage{multirow}
\usepackage{multicol}
\usepackage{booktabs}
\usepackage{float}
\usepackage{textcomp}
\usepackage{url}
\usepackage{hyperref}
\usepackage{siunitx}
\usepackage{xcolor}
\def\BibTeX{{\rm B\kern-.05em{\sc i\kern-.025em b}\kern-.08em
    T\kern-.1667em\lower.7ex\hbox{E}\kern-.125emX}}
\pdfobjcompresslevel=0
\usepackage{placeins}
\definecolor{MissingViolet}{RGB}{155,38,182}
\newif\ifshowdraftnotes
\showdraftnotesfalse
\newcommand{\regennote}[1]{\ifshowdraftnotes{{\color{MissingViolet}\textbf{[PLACEHOLDER: #1]}}}\fi}
\newcommand{\figupdate}[1]{\ifshowdraftnotes{{\color{red}\textbf{[FIGURE UPDATE BEFORE SUBMISSION: #1]}}}\fi}

\begin{document}

\title{Dynamics-Aware Federated Bayesian Optimization for Safe Controller Calibration in Heterogeneous Systems}

\author{Jakob Weber$^{1}$, Markus Gurtner$^{1}$, Adrian Trachte$^{2}$ and Andreas Kugi$^{3}$%
\thanks{$^{1}$Jakob Weber and Markus Gurtner are with the Center for Vision, Automation \& Control, AIT Austrian Institute of Technology GmbH, Vienna, Austria. {\tt\small jakob.weber@ait.ac.at}}
\thanks{$^{2}$Adrian Trachte is with Robert Bosch GmbH, Renningen, Germany.}
\thanks{$^{3}$Andreas Kugi is with the Automation and Control Institute, TU Wien, Austria, and the AIT Austrian Institute of Technology GmbH, Vienna, Austria.}}

\maketitle

\begin{abstract}
Accurate lateral path tracking across a heterogeneous vehicle fleet can require different feedforward steering actions, yet these differences need not align with fixed vehicle classes.
We study iterative federated learning of a residual feedforward term in a fixed two-degree-of-freedom steering controller, where the remaining feedback steering after each closed-loop learning stage forms the next local target.
FedAvg, proximal fine-tuning, IFCA, and FedEM are evaluated on a nonlinear three-platform fleet with 60 held-out test vehicles.
Across 10 matched training seeds, FedAvg reaches a fleet-mean lateral RMSE of $0.203\pm0.023\,\mathrm{m}$, while IFCA $K=4$ and FedEM $M=5$ reduce it to $0.172\pm0.013\,\mathrm{m}$ and $0.179\pm0.009\,\mathrm{m}$, respectively, with diminishing returns beyond about four models.
Despite this benefit, target-space diagnostics do not reveal persistent Compact/Sport/SUV groups.
Changing only the calibration trajectory increases RMSE by $6.2\%$ for IFCA $K=4$ and $1.8\%$ for FedEM $M=5$.
Under a common-trajectory intervention, the main variation in FedEM responsibility patterns correlates strongly with the understeer coefficient ($r=0.90$), while five-fold cross-validation gives $R^2=0.12$ from platform identity and $R^2=0.77$ from the understeer coefficient alone.
The results show that multi-model federated control can exploit continuous vehicle- and trajectory-dependent heterogeneity without recovering fixed vehicle classes.
\end{abstract}

\begin{IEEEkeywords}
Federated learning, feedforward control, vehicle dynamics, personalized federated learning, clustered federated learning.
\end{IEEEkeywords}

%%%%%%%%%%%%%%%%%%%%%%%%%%%%%%%%%%%%%%%%%%%%%%%%%%%%%%%%%%%%%%%%%%%%%%%%%%%%%%%%
\section{Introduction}
\label{sec:introduction}
%%%%%%%%%%%%%%%%%%%%%%%%%%%%%%%%%%%%%%%%%%%%%%%%%%%%%%%%%%%%%%%%%%%%%%%%%%%%%%%%
Accurate lateral path tracking is commonly realized by a two-degree-of-freedom (2DOF) steering controller in which feedback provides robustness and disturbance rejection while feedforward supplies the nominal steering action associated with the planned motion~\cite{kapania2015design}.
The feedforward term is difficult to calibrate across a heterogeneous fleet because physical characteristics like mass, yaw inertia, axle geometry, payload, and tire characteristics change the lateral-yaw dynamics.
A single feedforward map can therefore leave systematic vehicle-dependent corrections to the feedback controller.

A practical way to improve an imperfect feedforward term is to transfer repeatable feedback action into the feedforward path.
This idea is related to iterative learning control~\cite{arimoto1984bettering,bristow2006survey} and feedback-error learning~\cite{kawato1987hierarchical,miyamoto1988feedback}.
In this work, the objective is not a trial-specific command sequence.
Instead, each vehicle learns a parameterized map from local closed-loop rollouts, the updated map is redeployed, and new local data are generated for the next learning stage.
Federated learning (FL) enables this collaborative training for a fleet of vehicles without exchanging raw trajectories~\cite{mcmahan2017communication,kairouz2021advances,weber2024_survey}.
Figure~\ref{fig:framework} summarizes the resulting iterative control-learning loop and the federated model families considered in this work.

Previous work~\cite{weber2025_fed_FF} studied this feedforward-learning concept with a single FedAvg model for homogeneous vehicle dynamics.
The heterogeneous case introduces a different trade-off.
Pooling heterogeneous client data provides broad statistical support but can average vehicle-specific corrections.
Per-client fine-tuning reduces this bias but requires optimization on each new client.
Instead, the clustered federated learning (CFL) algorithm IFCA maintains several reusable models with hard client assignments~\cite{ghosh2020_IFCA}, while FedEM combines several reusable models through soft responsibilities~\cite{marfoq2021_FedEM}.
Most CFL studies focus on statistical heterogeneity in prediction tasks~\cite{el-rifai2025survey}.
In closed-loop control, however, the learning target also depends on the plant response and is regenerated after each controller update.

The central question is therefore broader than whether multi-model FL improves closed-loop tracking performance: \emph{what structure do useful federated control models represent?}
The results show that the answer is not a fixed vehicle-type specific partition.
Performance improves as IFCA and FedEM use more component models up to about four models, even though the vehicle-dependent learning targets do not form discrete groups that remain consistent across reference trajectories.
Controlled routing experiments instead show that the selected model or mixture depends on continuous vehicle dynamics and on the reference trajectory used for calibration.

Based on these observations, the contributions of this work are as follows:
\begin{itemize}
    \item We formulate iterative federated residual feedforward learning for a fixed 2DOF lateral steering controller. 
    Realized vehicle motion is available during training, while deployment remains reference driven.
    \item We compare FedAvg, proximal fine-tuning, and the two CFL algorithms IFCA and FedEM on a held-out heterogeneous vehicle fleet. 
    A sweep over two to five component models shows a robust benefit of model multiplicity, with diminishing returns beyond about four models.
    \item We analyze what the learned assignments represent. 
    Target-space diagnostics test whether discrete vehicle groups exist, IFCA support diagnostics quantify the cost of excessive hard partitioning, and controlled routing experiments show that the assignments depend on continuous dynamics descriptors and the calibration trajectory.
\end{itemize}
%
After introducing the iterative controller and the federated learning methods, Sec.~\ref{sec:performance} evaluates the closed-loop performance and the effect of model multiplicity.
Secs.~\ref{sec:conditions}--\ref{sec:expert_structure} then examine why additional models are useful and what structure the learned assignments represent.

\begin{figure}[t]
    \centering
    \includegraphics[width=\linewidth]{figures/graphical_abstract_2.png}
    \caption{Iterative federated learning of residual feedforward steering. 
    The remaining feedback action is transferred stage by stage into a pooled model, a set of hard-assigned component models (IFCA), a soft mixture of components (FedEM), or a client-personalized model (FedAvg-pFT) while raw rollout data remain local.} 
    \label{fig:framework}
\end{figure}

%%%%%%%%%%%%%%%%%%%%%%%%%%%%%%%%%%%%%%%%%%%%%%%%%%%%%%%%%%%%%%%%%%%%%%%%%%%%%%%%
\section{Iterative Federated Feedforward Learning}
\label{sec:method}
%%%%%%%%%%%%%%%%%%%%%%%%%%%%%%%%%%%%%%%%%%%%%%%%%%%%%%%%%%%%%%%%%%%%%%%%%%%%%%%%
\subsection{Vehicle model and 2DOF control architecture}
%
The plant of client $i$ is a nonlinear dynamic bicycle model with nonlinear tire forces, as used in~\cite{weber2025_fed_FF}.
Its state vector is $\boldsymbol{x}=[X,Y,\psi,v_x,v_y,r]^\top$, and its inputs are the applied front-axle steering angle $\delta_f$ and commanded longitudinal acceleration $a_x$.
The lateral-yaw dynamics are
%
\begin{align}
    \dot v_y &= \frac{F^{b}_{y,f}+F^{b}_{y,r}}{m}-r v_x, \\
    \dot r &= \frac{l_f F^{b}_{y,f}-l_r F^{b}_{y,r}}{I_z},
    \label{eq:lateral_dynamics}
\end{align}
%
where $v_x$ and $v_y$ denote the longitudinal and lateral velocities in the vehicle body frame, $r$ is the yaw rate, $m$ is the vehicle mass, and $I_z$ is the yaw moment of inertia. 
The distances from the center of gravity to the front and rear axles are denoted by $l_f$ and $l_r$, respectively. 
The lateral axle forces $F_{y,f}^b$ and $F_{y,r}^b$ are expressed in the vehicle body frame and are obtained from nonlinear Pacejka tire characteristics~\cite{pacejka2026magic}.
Combined longitudinal and lateral force demands are limited by a friction circle.
The remaining kinematics and longitudinal dynamics follow the standard nonlinear single-track formulation.
Longitudinal speed is tracked by the fixed longitudinal feedback loop of the baseline controller and is not learned in this study~\cite{althoff2014online}.

The lateral steering controller combines fixed feedback, nominal kinematic feedforward, and the learned residual,
%
\begin{align}
    \delta_{\rm FB}^{\rm raw} &= K_1 e_y + K_2 e_\psi + K_3 e_r, \\
    \delta_{\rm FF}^{\rm kin} &= \arctan\!\big((l_f+l_r)\kappa_{\rm ref}\big),\\
    \delta^{\rm cmd} &= \delta_{\rm FB}+\delta_{\rm FF}^{\rm kin}+f_{\boldsymbol\theta}(\boldsymbol z).
    \label{eq:2dof_compact}
\end{align}
%
Here, $e_y$, $e_\psi$, and $e_r$ denote the lateral-position, heading, and yaw-rate errors.
The raw feedback contribution is magnitude-limited before it is combined with the feedforward terms, and the resulting clipped contribution is denoted by $\delta_{\rm FB}$.
The combined steering command is subsequently subject to the steering-magnitude and steering-rate limits listed in Table~\ref{tab:fixed_settings} before it becomes the applied steering $\delta_f$~\cite{althoff2014online}.
Neither limit is active in the reported evaluation.
The learning target introduced below uses the pre-limit feedback demand $\delta_{\rm FB}^{\rm raw}$.

\subsection{Feedback-error target and outer iteration}
%
The learning process alternates between federated optimization and closed-loop data generation.
Within one outer stage, the client-local datasets are held fixed while the federated algorithm performs several communication rounds.
After the stage is complete, the updated server model or set of component models is deployed on all training vehicles to generate the local datasets for the next stage.
The initial \emph{Base} rollout uses no learned residual feedforward term.

Let $j\in\{0,\ldots,J-1\}$ denote the outer learning stage and define $f^{(-1)}\equiv0$.
At stage $j$, client $i$ follows its reference trajectory with the current feedforward model $f^{(j-1)}$ active.
The rollout produces the raw corrective feedback demand $\delta_{{\rm FB},i,k}^{\rm raw,(j)}$ and the local dataset
%
\begin{equation}
    \mathcal D_i^{(j)}=\left\{\left(\boldsymbol z_{i,k,\rm train}^{(j)},y_{i,k}^{(j)}\right)\right\}_{k=1}^{T}.
    \label{eq:stage_dataset}
\end{equation}
%
with $z_{k,\rm train}$ defined in~\eqref{eq:features}.
The target transfers the remaining systematic feedback correction into the feedforward path,
%
\begin{equation}
    y_{i,k}^{(j)} = f^{(j-1)}(\boldsymbol z_{k,\rm ref})     + \delta_{{\rm FB},i,k}^{\rm raw,(j)}.
    \label{eq:target}
\end{equation}
%
The stage-$j$ model is trained on $\mathcal D_i^{(j)}$ and is then used to generate the rollout for stage $j+1$.
If the learned feedforward term captures repeatable steering demand, the systematic raw feedback correction decreases across stages.

Training uses measured motion histories so that the network can learn from the realized vehicle response. 
During deployment, however, the learned term is evaluated only from reference signals, so that it remains a true feedforward contribution and does not rely on measured vehicle feedback.
With $\mathcal H_{N_h}(s_k)$ denoting the most recent $N_h$ samples of signal $s$, the training and deployment feature vectors are
%
\begin{align}
    \boldsymbol z_{k,\rm train} &=[\mathcal H_{N_h}(v_x),v_{x,\rm ref,k+r_p},\mathcal H_{N_h}(r),r_{\rm ref,k+r_p}]^\top, \label{eq:features} \\ 
    \boldsymbol z_{k,\rm ref} &=[\mathcal H_{N_h}(v_{x,\rm ref}),v_{x,\rm ref,k+r_p},\mathcal H_{N_h}(r_{\rm ref}),r_{\rm ref,k+r_p}]^\top \nonumber.
\end{align}
%
Vehicle-dependent response information is therefore available during learning, while the deployed neural term remains a feedforward function of the planned trajectory.

%%%%%%%%%%%%%%%%%%%%%%%%%%%%%%%%%%%%%%%%%%%%%%%%%%%%%%%%%%%%%%%%%%%%%%%%%%%%%%%%
\section{Federated Methods and Experimental Protocol}
\label{sec:experiment}
%%%%%%%%%%%%%%%%%%%%%%%%%%%%%%%%%%%%%%%%%%%%%%%%%%%%%%%%%%%%%%%%%%%%%%%%%%%%%%%%
\subsection{Federated model families}
%
FedAvg maintains one fleet-wide model~\cite{mcmahan2017communication}.
Selected clients start from the current server model, train it on their local stage dataset, and return only the update. The server forms a data-weighted average.

Instead, the hard clustered federated learning algorithm IFCA maintains $K$ component models~\cite{ghosh2020_IFCA}.
Each selected client evaluates the current models on local data and chooses
%
\begin{equation}
    c_i=\arg\min_{c\in\{1,\ldots,K\}} L_i(\boldsymbol\theta_c).
    \label{eq:ifca}
\end{equation}
%
The client trains only the selected model, and updates are aggregated separately for each component.
At test time, the same minimum-loss rule is applied to a model-independent calibration dataset.

The soft CFL algorithm FedEM maintains $M$ component models and uses soft responsibilities~\cite{marfoq2021_FedEM}.
For a test client, calibration losses define
%
\begin{align}
    \pi_{i,m}&=\frac{\exp[-L_i^{\rm cal}(\boldsymbol\theta_m)/\tau]}{\sum_{\ell=1}^{M}\exp[-L_i^{\rm cal}(\boldsymbol\theta_\ell)/\tau]},\\
    f_i(\boldsymbol z)&=\sum_{m=1}^{M}\pi_{i,m}f_{\boldsymbol\theta_m}(\boldsymbol z).
    \label{eq:fedem}
\end{align}
%
Thus, IFCA selects one reusable component model, whereas FedEM combines several component models.

FedAvg-pFT is included as a personalized federated learning baseline. 
Starting from the global FedAvg model, each test client adapts its parameters by minimizing the calibration loss with an $\ell_2$ penalty $\lambda_{\rm p}\|\boldsymbol\theta-\boldsymbol\theta_g\|_2^2/2$.
Unlike IFCA and FedEM, which select or combine reusable shared models, FedAvg-pFT performs gradient-based adaptation separately for each test client.

\subsection{Heterogeneous fleet}
%
The heterogeneous fleet is generated by sampling vehicle parameters around three nominal platforms: 
Compact, Sport, and SUV. 
Nominal vehicle parameters are listed in Table~\ref{tab:vehicle_parameters}, while the fixed controller, actuator, and measurement-noise settings are given in Table~\ref{tab:fixed_settings}. 
Platform labels are used only for fleet generation and subsequent analysis.

\begin{table}[t]
\centering
\caption{Nominal vehicle parameters. Client-to-client parameter distributions are given in the text.}
\label{tab:vehicle_parameters}
\scriptsize
\setlength{\tabcolsep}{2.3pt}
\begin{tabular}{@{}lccc@{}}
\toprule
Parameter & Compact & Sport & SUV\\
\midrule
$m$ [kg] & 1220 & 1520 & 2550\\
$I_z$ [kg m$^2$] & 1500 & 1800 & 5200\\
$l_f/l_r$ [m] & 1.05/1.45 & 1.15/1.35 & 1.30/1.55\\
$(B,C,E)_f$ & (9.0,1.27,0.93) & (13.5,1.38,0.86) & (7.5,1.22,0.95)\\
$(B,C,E)_r$ & (10.0,1.27,0.93) & (14.0,1.38,0.86) & (8.0,1.22,0.95)\\
\bottomrule
\end{tabular}
\end{table}

\begin{table}[t]
\centering
\caption{Fixed controller, actuator, and measurement-noise settings.}
\label{tab:fixed_settings}
\scriptsize
\setlength{\tabcolsep}{4pt}
\begin{tabular}{@{}ll@{}}
\toprule
Quantity & Value\\
\midrule
$K_1,K_2,K_3$ & $0.005,\ 0.1,\ 0.1$\\
Steering limit $\delta_{\max}$ & $0.209$ rad ($12^\circ$)\\
Rate limit $\dot\delta_{\max}$ & $5.236$ rad/s\\
Road friction $\mu$ & $0.9$\\
$\sigma_X,\sigma_Y$ & $0.02$ m\\
$\sigma_\psi$ & $0.001$ rad\\
$\sigma_{v_x}$ & $0.05$ m/s\\
$\sigma_r$ & $0.001$ rad/s\\
\bottomrule
\end{tabular}
\end{table}

Each client is generated from one of the three nominal platforms by introducing several sources of vehicle-to-vehicle variation.
First, the basic vehicle parameters are perturbed around their nominal values.
Mass $m$ and yaw inertia $I_z$ receive correlated log-normal perturbations with $\sigma_{\log}=0.03$ and correlation coefficient $\rho=0.85$, while the wheelbase is perturbed independently with $\sigma_{\log}=0.03$.
The Pacejka parameters $B$, $C$, and $E$ are varied using additive perturbations with standard deviations $0.05$, $0.015$, and $0.015$, respectively.
Next, a payload fraction $\gamma_{\rm pay}\sim\mathcal U(0,0.25)$ is placed at a longitudinal position $x_{\rm pay}/L\sim\mathcal U(-0.20,0.25)$.
The resulting payload is incorporated consistently into the vehicle mass, center-of-gravity position, axle distances, and yaw inertia.
Front and rear tire-condition factors are then sampled from $\mathcal U(0.85,1.10)$.
Finally, additional log-normal manufacturing scatter with $\sigma_{\log}=0.02$ is applied to $m$, $I_z$, $l_f$, $l_r$, $B_f$, and $B_r$.
No process noise is added.
The training fleet contains 150 vehicles, 50 per platform.

\subsection{Trajectories, training, and held-out evaluation}
%
The training set contains 50 reference trajectories.
Each trajectory is used by one Compact, one Sport, and one SUV vehicle, so the three platforms see the same trajectory set.
Longitudinal-speed commands are drawn from $\{20,25,30,33\}\,\mathrm{m/s}$, filtered with a first-order time constant of $3\,\mathrm{s}$, and limited to $|\dot v_x|\le2\,\mathrm{m/s^2}$.
The lateral generator combines straight and mildly curved road segments with single- and double-lane changes.
Event curvature is limited such that $|v_{\rm ref}^2\kappa_{\rm ref}|\le7\,\mathrm{m/s^2}$.
Rollouts last $25\,\mathrm{s}$ with $dt=0.05\,\mathrm{s}$ and RK4 integration.

All methods use the same $[12,5,5,5,1]$ ReLU network, trained with Adam at a learning rate of $10^{-3}$.
The feature vector contains the $N_h=5$ most recent samples of $v_x$ and $r$, together with a $0.30\,\mathrm{s}$ reference preview corresponding to $r_p=6$ samples.
Training uses six outer stages.
In each communication round, 15 of the 150 training clients are sampled.
Each sampled client receives the current server model or component set, optimizes its assigned local model on the current rollout dataset for $E=3$ local epochs, and returns only the model update.
After 20 communication rounds, the outer stage ends and the resulting server model or component set is redeployed on all 150 training vehicles to generate the next local datasets.
FedEM uses soft responsibilities with temperature annealing from $\tau=1$ to $0.05$.
FedAvg-pFT uses 10 adaptation epochs, learning rate $3\times10^{-4}$, and $\lambda_{\rm p}=0.01$.
The main performance results are evaluated over 10 matched training seeds.
The model-count sweep considers IFCA $K\in\{2,3,4,5\}$ and FedEM $M\in\{2,3,4,5\}$.

After the architecture and training protocol are fixed, a held-out test fleet of 60 vehicles, 20 per platform, and 20 new matched reference trajectories is generated.
For client $i$, tracking performance is measured by the RMSE of the lateral error $J_i=\sqrt{T^{-1}\sum_{k=1}^{T} e_{y,i,k}^{2}}$.
Table~\ref{tab:main_results} also reports the mean peak error $N^{-1}\sum_i\max_k|e_{y,i,k}|$, which captures the average worst lateral deviation within a rollout.

For a previously unseen test client, IFCA, FedEM, and pFT require local calibration data to determine the model assignment, mixture responsibilities, or client-specific adaptation, respectively.
We therefore use one additional model-independent calibration rollout with the learned residual disabled, i.e., with only the fixed feedback and kinematic feedforward terms active.
Calibration and evaluation use the same reference trajectory with independent measurement noise unless Sec.~\ref{sec:expert_structure} explicitly changes the calibration trajectory.

% Place the full-width performance table before Sec. IV so IEEEtran can float it
% to the top of the page on which the results discussion begins.
\input{figures/table_main_results}

%%%%%%%%%%%%%%%%%%%%%%%%%%%%%%%%%%%%%%%%%%%%%%%%%%%%%%%%%%%%%%%%%%%%%%%%%%%%%%%%
\section{Closed-Loop Performance and Model Multiplicity}
\label{sec:performance}
%%%%%%%%%%%%%%%%%%%%%%%%%%%%%%%%%%%%%%%%%%%%%%%%%%%%%%%%%%%%%%%%%%%%%%%%%%%%%%%%

Table~\ref{tab:main_results} shows that learned feedforward control improves both mean and peak tracking error.
The fixed feedback plus kinematic-feedforward controller gives $0.245\,\mathrm{m}$ fleet RMSE.
FedAvg reduces this to $0.203\pm0.023\,\mathrm{m}$, while FedAvg-pFT reaches $0.191\pm0.023\,\mathrm{m}$.
The lowest IFCA mean is $0.172\,\mathrm{m}$ for both $K=4$ and $K=5$, and FedEM reaches $0.179\pm0.009\,\mathrm{m}$ at $M=5$.
Relative to FedAvg, these are reductions of about $15.3\%$ for IFCA $K=4/5$ and $11.8\%$ for FedEM $M=5$.
All three platforms improve at the higher model counts.
Mean peak error also decreases from $0.556\,\mathrm{m}$ for FedAvg to $0.493/0.492\,\mathrm{m}$ for IFCA $K=4/5$ and $0.505\,\mathrm{m}$ for FedEM $M=5$.

Figure~\ref{fig:stage_mechanism} shows how the improvement develops across the outer stages.
Panel (a) uses one held-out SUV trajectory selected solely from the reference trajectory: 
its RMS curvature-rate intensity is closest to the median of the 20 test trajectories, without using any method's tracking error.
Panels (b)--(c) show averages over all 10 training seeds.
Most of the reduction occurs during the first two to three stages.
The fleet RMSE of IFCA $K=4$ and FedEM $M=5$ then approaches a low plateau below FedAvg.
The raw feedback demand shows the same transfer of systematic correction into feedforward control. 
IFCA reaches its minimum around stage 2 and then fluctuates slightly, while FedEM decreases more gradually through the later stages.

\begin{figure*}[t]
    \centering
    \begin{subfigure}[t]{0.325\textwidth}
        \centering
        \includegraphics[width=\linewidth]{figures/fig2_a_time_trace.pdf}
        \caption{Representative $e_y(t)$ trace.}
    \end{subfigure}
    \hfill
    \begin{subfigure}[t]{0.325\textwidth}
        \centering
        \includegraphics[width=\linewidth]{figures/fig2_b_round_RMSE.pdf}
        \caption{Fleet-mean RMSE.}
    \end{subfigure}
    \hfill
    \begin{subfigure}[t]{0.325\textwidth}
        \centering
        \includegraphics[width=\linewidth]{figures/fig2_c_round_FB.pdf}
        \caption{Raw feedback demand.}
    \end{subfigure}
    \caption{Closed-loop and stage-wise behavior on the held-out test fleet. 
    Panel (a) shows one representative trajectory selected from reference-trajectory transient intensity without inspecting method errors. 
    Panels (b)--(c) show mean $\pm$ sample standard deviation over 10 matched training seeds for IFCA $K=4$ and FedEM $M=5$. 
    Base denotes the controller before neural feedforward learning, and stage 0 is the first trained model.}
    \label{fig:stage_mechanism}
\end{figure*}

Figure~\ref{fig:model_multiplicity} shows the main model-count result.
Both IFCA and FedEM improve as the number of component models increases from two to about four, followed by diminishing returns.
The best-performing counts are therefore four to five models, although the fleet was generated from only three nominal platform families.
The error bars show across-seed standard deviation, while statistical comparisons use paired differences between methods trained with the same seed.
Relative to FedAvg, all multi-model settings except IFCA $K=2$ show a significant paired improvement around $5\%$.
Relative to the stronger FedAvg-pFT baseline, significant paired improvements are obtained only for IFCA $K=4$ ($p=0.008$), IFCA $K=5$ ($p=0.025$), and FedEM $M=5$ ($p=0.047$).
Thus, low model counts do not reliably outperform client-specific fine-tuning, while the benefit of higher multiplicity is seed-consistent.

\begin{figure}[t]
    \centering
    \includegraphics[width=0.96\linewidth]{figures/fig3_model_multiplicity.pdf}
    \caption{Effect of model multiplicity on fleet-mean RMSE over 10 matched training seeds. 
    Markers and error bars show mean $\pm$ sample standard deviation. 
    The FedAvg-pFT line shows its across-seed mean. 
    Paired comparisons use matched training seeds and are reported in the text.}
    \label{fig:model_multiplicity}
\end{figure}
%
The plateau also has an assignment-support interpretation for hard clustering.
Table~\ref{tab:ifca_support} shows stable final-stage assignments through $K=4$.
At $K=5$, two of 10 runs contain an empty cluster and final-stage reassignment rises to $59.1\pm23.6\%$.
The held-out RMSE does not deteriorate, but the fifth component gives no further mean improvement while the assignment becomes much less stable.
Hence, $K=4$ provides the clearest trade-off between tracking performance and stable hard assignments in this experiment, without implying a unique RMSE optimum.

\input{figures/table_ifca_support.tex}

%%%%%%%%%%%%%%%%%%%%%%%%%%%%%%%%%%%%%%%%%%%%%%%%%%%%%%%%%%%%%%%%%%%%%%%%%%%%%%%%
\section{Do Useful Component Models Correspond to Natural Vehicle Groups?}
\label{sec:conditions}
%%%%%%%%%%%%%%%%%%%%%%%%%%%%%%%%%%%%%%%%%%%%%%%%%%%%%%%%%%%%%%%%%%%%%%%%%%%%%%%%
The multiplicity result creates an important question.
More models improve control even though the fleet was generated from only three nominal platforms.
We therefore test whether the useful component models recover discrete vehicle groups that remain consistent across reference trajectories, or whether they approximate a more continuous control problem.

\subsection{Shared and vehicle-dependent control targets}
%
For each training trajectory $q$, let $\mathcal I_q$ denote the three vehicles that follow this same reference.
We decompose the stage-zero feedback-error target into the trajectory-shared signal
$\bar y_q(t)$ and the vehicle-dependent deviation $d_{i,q}(t)$,
%
\begin{align}
    \bar y_q(t)&=\frac{1}{|\mathcal I_q|}\sum_{i\in\mathcal I_q}y_{i,q}(t), \nonumber \\
    d_{i,q}(t)&=y_{i,q}(t)-\bar y_q(t).
\label{eq:target_decomposition}
\end{align}
%
Thus, $\bar y_q$ is the correction common to the three vehicles on trajectory $q$, while $d_{i,q}$ contains the additional correction associated with vehicle $i$.
Let $S_{\rm shared}(\nu)$ denote the spectral energy density of $\bar y_q$, averaged over trajectories, and let $S_{\rm dyn}(\nu)$ denote the corresponding spectral energy density of $d_{i,q}$, averaged over vehicles and trajectories.
We use $\eta_{\rm dyn}(\mathcal B)$ to measure which fraction of the target energy in frequency band $\mathcal B$ is caused by vehicle-to-vehicle differences rather than by the correction shared across the matched vehicles,
%
\begin{equation}
\eta_{\rm dyn}(\mathcal B)=
\frac{\int_{\mathcal B}S_{\rm dyn}(\nu)\,d\nu}
{\int_{\mathcal B}[S_{\rm shared}(\nu)+S_{\rm dyn}(\nu)]\,d\nu}.
\label{eq:relative_target_heterogeneity}
\end{equation}
%
The reference trajectories contain about $25\%$, $45\%$, and $19\%$ of their yaw-rate-reference energy in $0$--$0.15$, $0.15$--$0.5$, and $0.5$--$2\,\mathrm{Hz}$, respectively.
Thus, the dataset excites the frequency range in which transient lateral dynamics matter.
Figure~\ref{fig:target_structure}(a) shows $\eta_{\rm dyn}=0.48$, $0.32$, and $0.09$ over the same bands.
In particular, even though $19\%$ of the reference energy lies in the $0.5$--$2\,\mathrm{Hz}$ band, only about $9\%$ of the corresponding target energy is vehicle dependent.
The learning problem therefore contains genuine vehicle-specific corrections, but much of the target remains shared across the fleet.

\subsection{Target-space clusterability}
%
We next ask whether the vehicle-dependent target deviations form discrete groups.
For vehicle $i$ on trajectory $q$, $d_{i,q}(t)$ is a time signal with $T$ rollout samples.
We stack these samples into the target-deviation vector
%
\begin{equation}
    \boldsymbol d_{i,q}=[d_{i,q}(1),\ldots,d_{i,q}(T)]^\top\in\mathbb R^T.
\label{eq:target_signature}
\end{equation}
%
PCA is applied to the target-deviation vectors, and the first 20 principal components are retained for K-means clustering.
Figure~\ref{fig:target_structure}(b) shows the PC1--PC2 projection, in which the Compact, Sport, and SUV samples overlap strongly.
Cluster separation is quantified by the silhouette score, where values near one indicate well-separated groups and values near zero indicate substantial overlap~\cite{rousseeuw1987silhouettes}.
After excluding one trajectory-level outlier shared across all three platforms, the silhouette scores remain low for $K=2,\ldots,5$, ranging from $0.166$ to $0.204$.

A persistent vehicle grouping could still be hidden by trajectory-specific variation.
To test this, target-deviation vectors from the same physical vehicle are averaged over $N_{\rm traj}=1,2,$ or $4$ independent reference trajectories before repeating the PCA/K-means analysis.
If trajectory-specific effects were masking stable vehicle groups, this averaging should increase the silhouette score.
Figure~\ref{fig:target_structure}(c) shows the opposite: the $K=3$ silhouette decreases from $0.166$ to $0.147$ and $0.068$ as $N_{\rm traj}$ increases.
Averaging trajectory-specific variation therefore does not reveal a hidden Compact/Sport/SUV partition.

\begin{figure*}[t]
    \centering
    \begin{subfigure}[t]{0.27\textwidth}
        \centering
        \includegraphics[width=\linewidth]{figures/fig4a_eta_dyn.pdf}
        \caption{Relative target heterogeneity.}
    \end{subfigure}
    \hfill
    \begin{subfigure}[t]{0.34\textwidth}
        \centering
        \includegraphics[width=\linewidth]{figures/fig4b_target_pca.pdf}
        \caption{Target-space PCA projection.}
    \end{subfigure}
    \hfill
    \begin{subfigure}[t]{0.31\textwidth}
        \centering
        \includegraphics[width=\linewidth]{figures/fig4c_silhouette.pdf}
        \caption{Trajectory averaging.}
    \end{subfigure}
    \caption{Structure of the vehicle-dependent learning target, computed from the training population. 
    Panel (a) shows the fraction of target energy attributed to vehicle-dependent deviations in three frequency bands. 
    Panle (b) depcits the PC1--PC2 projection of the target-deviation vectors. 
    The clustering analysis itself uses 20 retained principal components. 
    Panel (c) shows the $K=3$ silhouette score after averaging target-deviation vectors over multiple independent trajectories of the same physical vehicles. 
    The absence of a stable platform-aligned grouping contrasts with the performance benefit of using several component models in Fig.~\ref{fig:model_multiplicity}.}
    \label{fig:target_structure}
\end{figure*}

These results show that useful specialization does not require naturally separated vehicle groups. 
Instead, IFCA can use a finite set of component models to approximate different regions of a continuously varying control mapping, while FedEM can represent related structure through soft mixtures of shared models. 
The remaining question is therefore not whether the learned assignments recover physical vehicle classes, but which vehicle and reference-trajectory properties determine how clients are assigned to the available models. Section~\ref{sec:expert_structure} addresses this question.

%%%%%%%%%%%%%%%%%%%%%%%%%%%%%%%%%%%%%%%%%%%%%%%%%%%%%%%%%%%%%%%%%%%%%%%%%%%%%%%%
\section{What Do the Learned Assignments Represent?}
\label{sec:expert_structure}
%%%%%%%%%%%%%%%%%%%%%%%%%%%%%%%%%%%%%%%%%%%%%%%%%%%%%%%%%%%%%%%%%%%%%%%%%%%%%%%%
Let $\boldsymbol\xi_i$ denote, in an abstract sense, the continuous physical properties that determine the lateral dynamics of vehicle $i$.
Examples used below are the dominant yaw-response time, the understeer coefficient, and steady-state gains.
Let $\boldsymbol\chi_q$ denote features of reference trajectory $q$, including its speed and curvature profile.
We use two controlled interventions to test whether the learned assignments depend on the calibration trajectory, on vehicle dynamics, or on both.

\subsection{The calibration trajectory affects routing}
%
In the held-out evaluation of Table~\ref{tab:main_results} and Fig.~\ref{fig:model_multiplicity}, IFCA and FedEM infer their model assignment or mixture from a calibration rollout that uses the same reference trajectory as the subsequent evaluation, with independent measurement noise.
To isolate the role of that trajectory, we keep the physical vehicle, trained models, and evaluation trajectory fixed and change only the reference trajectory used for calibration.
Figure~\ref{fig:routing_structure}(a) evaluates IFCA $K=4$ and FedEM $M=5$, which are representative high-performing configurations in Table~\ref{tab:main_results}, and reports fleet means over the three mechanism-analysis seeds.
FedEM changes from $0.1745\pm0.0018$ to $0.1777\pm0.0036\,\mathrm{m}$, a $1.8\%$ degradation, while remaining below its uniform-mixture reference of $0.1908\,\mathrm{m}$.
IFCA changes from $0.1625\pm0.0031$ to $0.1726\pm0.0034\,\mathrm{m}$, a $6.2\%$ degradation.
The selected IFCA component and the FedEM mixture therefore depend partly on the reference trajectory used for calibration, with hard IFCA routing showing the larger sensitivity in this experiment.
% TODO before submission: if the cross-trajectory evaluation remains inexpensive, extend Fig. 5(a) from 3 to 10 matched seeds for consistency with the main performance study.

\begin{figure*}[t]
    \centering
    \begin{subfigure}[t]{0.35\textwidth}
        \centering
        \includegraphics[width=\linewidth]{figures/fig5a_context_routing.pdf}
        \caption{Sensitivity to calibration trajectory.}
    \end{subfigure}
    \hfill
    \begin{subfigure}[t]{0.34\textwidth}
        \centering
        \includegraphics[width=\linewidth]{figures/fig5b_routing_kus_scatter.pdf}
        \caption{Dominant routing coordinate versus $K_{\rm us}$.}
    \end{subfigure}
    \hfill
    \begin{subfigure}[t]{0.27\textwidth}
        \centering
        \includegraphics[width=\linewidth]{figures/fig5c_routing_cv.pdf}
        \caption{Cross-validated prediction of $z_1$.}
    \end{subfigure}
    \caption{The learned assignments depend on both the calibration trajectory and continuous vehicle dynamics. 
    Panel (a) shows that changing only the calibration trajectory increases fleet RMSE more strongly for IFCA $K=4$ than for FedEM $M=5$. 
    Open and filled markers show same- and cross-trajectory calibration. 
    Panel (b) depcits the results from the common-trajectory FedEM $M=5$ experiment, in which the dominant responsibility-similarity coordinate $z_1$ is strongly related to the understeer coefficient. 
    Marker shape denotes nominal platform. 
    Panel (c) shows the five-fold cross-validated coefficient of determination for predicting the fixed coordinate $z_1$. Continuous dynamics descriptors, especially $K_{\rm us}$, predict the learned routing coordinate much better than the nominal platform label.}
    \label{fig:routing_structure}
\end{figure*}

\subsection{Continuous dynamics organize the dominant FedEM routing coordinate}
%
To isolate the vehicle-dynamics contribution, we use a controlled common-trajectory experiment in which all vehicles share the same reference trajectory and FedEM is retrained under this distribution.
We analyze $M=5$, the best-performing FedEM configuration in Table~\ref{tab:main_results}.
For each of the three independently trained seeds, every held-out vehicle $i$ has a responsibility vector $\boldsymbol\pi_i^{(s)}\in\mathbb R^5$ that describes how strongly the five component models contribute to its prediction.
Directly averaging these vectors across seeds is not meaningful because component labels may be permuted between independent training runs.

Instead, we compare clients through the similarity of their responsibility vectors.
For seed $s$, define
%
\begin{equation}
A_{ij}^{(s)}=\boldsymbol\pi_i^{(s)\top}\boldsymbol\pi_j^{(s)},
\qquad
\bar A=\frac{1}{S}\sum_{s=1}^{S}A^{(s)}.
\label{eq:routing_affinity}
\end{equation}
%
Here, $A^{(s)}$ is the client-similarity matrix for seed $s$, and $\bar A$ is the corresponding similarity matrix averaged over the $S$ independent seeds.
The entry $A_{ij}^{(s)}$ is large when clients $i$ and $j$ place responsibility on the same components.
Moreover, the dot product is unchanged if the component indices are permuted in the same way for all clients within a seed.
With 60 held-out vehicles, $A^{(s)},\bar A\in\mathbb{R}^{60\times60}$.
Kernel PCA is applied to $\bar A$ to obtain a low-dimensional representation of this learned routing similarity.
The leading coordinate $z_1$ corresponds to $66\%$ of the summed positive kernel-PCA eigenvalues and is used below as the dominant routing coordinate.

We compare $z_1$ with four continuous dynamics descriptors computed from the linearized two-state bicycle model at the mean speed of the common trajectory.
Let $A_i$ denote the standard linearized $[v_y,r]^\top$ bicycle-model state matrix for vehicle $i$ at this speed.
The considered descriptors are the dominant-pole yaw-response time $\tau_r$, the understeer coefficient $K_{\rm us}$, and the steady-state yaw-rate and sideslip gains.
The understeer coefficient is
%
\begin{equation}
    K_{\rm us}=\frac{m}{L}\left(\frac{l_r}{C_{\alpha,f}}-\frac{l_f}{C_{\alpha,r}}\right),
\label{eq:kus}
\end{equation}
%
where $C_{\alpha,f}$ and $C_{\alpha,r}$ are the zero-slip cornering stiffnesses obtained from the local Pacejka slopes.
% The remaining descriptors are the steady-state yaw-rate gain and sideslip gain.

Figure~\ref{fig:routing_structure}(b) compares the dominant kernel-PCA coordinate $z_1$, obtained from the routing-similarity matrix $\bar A$, with the understeer coefficient $K_{\rm us}$.
Since the sign of a kernel-PCA coordinate is arbitrary, $z_1$ is oriented such that its Pearson correlation with $K_{\rm us}$ is positive.
The resulting correlation of $r(z_1,K_{\rm us})=0.90$ shows that the dominant variation in the FedEM responsibility patterns is strongly associated with this continuous vehicle-dynamics descriptor.
At the same time, the platform markers overlap along $z_1$ rather than forming three distinct Compact, Sport, and SUV groups.
This suggests that the observed organization is related to continuous vehicle dynamics rather than to nominal platform identity alone.

Figure~\ref{fig:routing_structure}(c) tests this interpretation quantitatively.
For each considered predictor set---platform identity, $\tau_r$, $K_{\rm us}$, or all four continuous dynamics descriptors jointly---a linear model is fitted to predict the fixed routing coordinate $z_1$ and evaluated using five-fold cross-validation.
The resulting coefficient of determination is reported as $R^2$.
The platform identity alone explains little of the variation in $z_1$ ($R^2=0.12$), whereas the continuous dynamics descriptors provide substantially stronger predictions.
In particular, $\tau_r$ gives $R^2=0.29$, $K_{\rm us}$ alone gives $R^2=0.77$, and all four descriptors jointly give $R^2=0.80$.
Adding the platform label to the four descriptors does not improve the prediction ($R^2=0.79$, not shown).
As an additional control, subtracting the corresponding platform-wise means from both $z_1$ and $K_{\rm us}$ changes their Pearson correlation only from $0.891$ to $0.889$.
Hence, the strong relation between $z_1$ and $K_{\rm us}$ persists within the nominal platform groups and cannot be explained simply by differences between Compact, Sport, and SUV vehicles.
This analysis concerns only the dominant routing coordinate and does not imply that the complete five-component responsibility vector is determined by a single physical quantity.

Together, the two interventions support a task-dependent interpretation of the multi-model controllers.
With $\boldsymbol\xi_i$ denoting continuous vehicle-dynamics properties and $\boldsymbol\chi_q$ the features of the reference trajectory used for calibration, IFCA can be represented conceptually as
%
\begin{equation}
    f_{i,q}(\boldsymbol z)\approx f_{c(\boldsymbol\xi_i,\boldsymbol\chi_q)}(\boldsymbol z),
\label{eq:ifca_joint_structure}
\end{equation}
%
while FedEM uses
%
\begin{equation}
    f_{i,q}(\boldsymbol z)\approx
\sum_{m=1}^{M}\pi_m(\boldsymbol\xi_i,\boldsymbol\chi_q)f_m(\boldsymbol z).
\label{eq:fedem_joint_structure}
\end{equation}
%
The assignments are therefore better interpreted as task-dependent coordinates of heterogeneous control mappings than as fixed vehicle-class labels.

\FloatBarrier

%%%%%%%%%%%%%%%%%%%%%%%%%%%%%%%%%%%%%%%%%%%%%%%%%%%%%%%%%%%%%%%%%%%%%%%%%%%%%%%%
\section{Discussion}
\label{sec:discussion}
%%%%%%%%%%%%%%%%%%%%%%%%%%%%%%%%%%%%%%%%%%%%%%%%%%%%%%%%%%%%%%%%%%%%%%%%%%%%%%%%
The results change the usual interpretation of clustered FL in this control problem.
Physical heterogeneity can make different control models beneficial for different vehicles, but useful multi-model control does not require pre-existing discrete client groups.
The best-performing model counts are four to five although the fleet has three nominal platforms, and the target-space analysis does not recover those platforms.
The performance gain is therefore more consistent with a finite set of models approximating continuously varying control mappings than with discovery of hidden vehicle classes.

Hard and soft multi-model methods have different practical properties.
IFCA attains the lowest mean RMSE in the present sweep, but $K=5$ produces unstable assignments and $K=4$ is more sensitive than FedEM $M=5$ to a change in calibration trajectory.
FedEM improves more gradually with model count and its soft mixture is less sensitive to this mismatch.
We therefore do not infer a universal ranking between the two methods. The appropriate choice depends on client support, the calibration protocol, and the desired robustness of model assignment.

The study has several limitations.
It is simulation based and uses one nonlinear single-track/Pacejka model family, one fixed feedback-controller structure, and one class of highway-style reference trajectories.
We do not provide a nonlinear convergence or stability proof for the outer learning loop.
IFCA and FedEM require a calibration rollout whose representativeness affects performance, while the common-trajectory experiment is intended only as a controlled diagnostic.
Future work should study online routing from estimated physical descriptors, extend the analysis to richer vehicle and operating-condition variation, and validate the method on experimental fleet data.

%%%%%%%%%%%%%%%%%%%%%%%%%%%%%%%%%%%%%%%%%%%%%%%%%%%%%%%%%%%%%%%%%%%%%%%%%%%%%%%%
\section{Conclusion}
\label{sec:conclusion}
%%%%%%%%%%%%%%%%%%%%%%%%%%%%%%%%%%%%%%%%%%%%%%%%%%%%%%%%%%%%%%%%%%%%%%%%%%%%%%%%
This work examined how federated feedforward steering controllers can represent heterogeneous vehicle lateral dynamics without exchanging raw vehicle trajectories.
On a held-out heterogeneous three-platform test fleet, the fixed feedback plus kinematic-feedforward controller gives $0.245\,\mathrm{m}$ fleet RMSE and FedAvg reduces it to $0.203\,\mathrm{m}$.
Multi-model IFCA and FedEM reduce the error further as model multiplicity increases to about four components.
IFCA $K=4$ and $K=5$ both reach $0.172\,\mathrm{m}$, but $K=4$ provides the better assignment-stability trade-off, while FedEM $M=5$ reaches $0.179\,\mathrm{m}$.

The main result is that these gains do not require recovery of fixed vehicle classes.
The vehicle-dependent control targets do not form stable Compact/Sport/SUV groups across reference trajectories.
The learned assignments depend on the calibration trajectory, and FedEM clients with similar responsibility patterns are organized much more strongly by continuous dynamics descriptors, such as the understeer coefficient $K_{\rm us}$, than by platform identity.
Multi-model federated control should therefore be interpreted as a task-driven approximation of heterogeneous control mappings rather than as vehicle-class discovery.

\bibliographystyle{IEEEtran}
\bibliography{references}

\end{document}
